\documentclass[11pt]{article}
\usepackage[T1]{fontenc}
\usepackage{lmodern,microtype,amsmath,amssymb,amsthm}
\usepackage[margin=1in]{geometry}
\usepackage[colorlinks=true,allcolors=blue]{hyperref}
\newtheorem{theorem}{Theorem}
\newtheorem{lemma}[theorem]{Lemma}
\newtheorem{corollary}[theorem]{Corollary}
\newcommand{\norm}[1]{\left\lVert#1\right\rVert_\infty}
\newcommand{\Z}{\mathbb Z}
\title{Fixed-Width Gaussian Resampling Is Polynomially Conditioned}
\author{Research note prepared with OpenAI Codex}
\date{October 13, 2026}
\begin{document}
\maketitle
\begin{abstract}
An exact bilateral Newton factorization proves that the Gaussian
resampling inverse has norm $O(\theta^{-2})$ at a fixed Gaussian width,
where $\theta>0$ is the oversampling gap. This removes the width-growth
requirement from the conditioning argument. The proof uses convergent
infinite triangular factors and then periodizes them; it does not assert
operator-norm convergence of finite sections. A fast implementation and
its tensor cost remain separate obligations. No improved final
integer-multiplication exponent is claimed in this note.
\end{abstract}

Let $1\le s<t\le2s$ be integers, $\rho=t/s=1+\theta$, and $a\ge1$.
For every $i\in\Z$ set
\[
 q_i=\lfloor\rho i+1/2\rfloor,\quad
 \beta_i=\rho i-q_i\in[-1/2,1/2),\quad
 x_i=e^{2a\rho q_i},\quad
 v_i=e^{-a\rho^2i^2+a\beta_i^2}.
\]
Thus $q_{i+s}=q_i+t$ and $\beta_{i+s}=\beta_i$.
The bilateral unwrapped matrix and the finite periodic matrix are
\[
 G_{ij}=e^{-a[(\rho j-q_i)^2-\beta_j^2]},\qquad
 N_{ij}=\sum_{n\in\Z}G_{i,j+ns}\quad(0\le i,j<s).
\]
Norms of finite matrices are maximum absolute row sums. Bilateral
matrices act on $\ell^\infty(\Z)$ and have the same absolute-row-sum norm.

\begin{theorem}\label{thm:main}
Put
\[
 c=\prod_{m=1}^{\infty}(1-e^{-2a\rho m}).
\]
Both $G$ and $N$ are invertible bounded operators, and
\[
 \norm{G^{-1}},\ \norm{N^{-1}}
 \le \frac{c^{-4}e^{a/2}}{(1-e^{-a\rho\theta})^2}.
\]
In particular, for each fixed $a\ge1$ this is $O_a(\theta^{-2})$.
At $a=1$ the explicit bound $\norm{N^{-1}}<16\theta^{-2}$ holds.
There is no coprimality condition.
\end{theorem}

For a finite or summable positive alphabet $X$, write $e_h(X)$ and
$h_h(X)$ for its elementary and complete homogeneous symmetric
polynomials, with $e_0=h_0=1$.

\begin{lemma}[The partition bounds]\label{lem:partition}
If successive positive nodes increase by a factor at least
$R=e^{2a\rho}$, then the following quantities are at most $c^{-1}$:
\[
 \frac{e_h(x_l,\ldots,x_{j-1})}
      {x_{j-h}\cdots x_{j-1}},\qquad
 \frac{h_h(x_l,\ldots,x_i)}{x_i^h}.
\]
The same bounds hold with $l=-\infty$.
\end{lemma}
\begin{proof}
In the elementary polynomial compare each increasing $h$-subset with
the $h$ largest indices. Its deficits form a partition, and its relative
weight is at most $R$ to the negative partition size. In the complete
homogeneous polynomial, a multiset compared with $h$ copies of the
largest node has the same partition bound. Summing over all partitions
gives $\prod_{m\ge1}(1-R^{-m})^{-1}=c^{-1}$.
The infinite alphabets converge because $x_k/x_i\le R^{k-i}$ for
$k\le i$; monotone limits preserve the inequalities.
\end{proof}

\begin{proof}[Proof of Theorem~\ref{thm:main}]
First restrict $G$ to an arbitrary finite interval $[l,r]$.
The usual Newton expansion of monomials gives an exact factorization
\[
 G_{[l,r]}=L^{(l)}D^{(l)}U^{(l)},\qquad
 D^{(l)}_{ii}=C_i^{(l)}:=\prod_{l\le k<i}(1-x_k/x_i),
\]
where both triangular factors have diagonal one and
\begin{align*}
 L^{(l)}_{ij}
 &=G_{ij}\frac{\prod_{l\le k<j}(1-x_k/x_i)}{C_j^{(l)}}
 &&(i\ge j),\\
 U^{(l)}_{ij}
 &=\frac{v_j}{v_i}h_{j-i}(x_l,\ldots,x_i)
 &&(j\ge i).
\end{align*}
All other entries are zero. These formulas also fix the normalization
when $l$ is negative: shift the index origin to $l$, use nodes
$x'_k=e^{-2a\rho^2l}x_k$ and column factors
$v'_j=e^{2a\rho^2lj-a\rho^2l^2}v_j$; the extra factors cancel.

Since $q_{k+1}-q_k\ge1$, we have $c\le C_i^{(l)}\le1$.
Lemma~\ref{lem:partition} and
$(v_j/v_i)x_i^{j-i}=G_{ij}$ show entrywise that
\[
 0\le L^{(l)}_{ij}\le c^{-1}G_{ij},\qquad
 0\le U^{(l)}_{ij}\le c^{-1}G_{ij}.
\]
The Gaussian envelope
\[
 G_{i,i+h}\le e^{-a(\rho^2h^2-\rho|h|)}
\]
is summable uniformly in $i$. Every factor entry has a limit as
$l\to-\infty$. Products of the limiting factors are therefore justified
by dominated convergence, and give bounded triangular matrices $L,U$
and diagonal $D$ with $G=LDU$.

For completeness, the inverse entries also have explicit finite-origin
formulas. If $h=i-j\ge0$, divided differences give
\[
 (L^{(l)})^{-1}_{ij}
 =(-1)^h\frac{C_i^{(l)}}{C_j^{(l)}}
 \frac{\exp\!\left(a[-\beta_i^2+\beta_j^2-\rho^2h
              +2\rho\sum_{k=j+1}^{i}\beta_k]\right)}
      {\prod_{k=j+1}^{i}(1-x_j/x_k)}.
\]
If $h=j-i\ge0$, the inverse Newton expansion gives
\[
 (U^{(l)})^{-1}_{ij}
 =(-1)^h\frac{v_j}{v_i}e_h(x_l,\ldots,x_{j-1}).
\]
Using $|\beta_k|\le1/2$ and the partition bound yields, uniformly in $l$,
\[
 |(L^{(l)})^{-1}_{ij}|
 \le c^{-2}e^{a/4}e^{-a\rho\theta(i-j)},\qquad
 |(U^{(l)})^{-1}_{ij}|
 \le c^{-1}e^{a/4}e^{-a\rho\theta(j-i)}.
\]
For the second estimate the dominant exponent is exactly
\[
 a\left[\beta_j^2-\beta_i^2-\rho^2h
               -2\rho\sum_{k=i}^{j-1}\beta_k\right].
\]
These inverse entries converge as $l\to-\infty$. Each entry of a
product of two lower triangular factors, or of two upper triangular
factors, involves only finitely many intermediate indices. Hence the
limiting inverse identities hold entrywise. The uniform summable row
bounds promote them to identities of bounded operators on
$\ell^\infty(\Z)$. Consequently
\[
 G^{-1}=U^{-1}D^{-1}L^{-1},\qquad
 \norm{G^{-1}}
 \le\frac{c^{-4}e^{a/2}}{(1-e^{-a\rho\theta})^2}.
\]

Every limiting factor and inverse is invariant under simultaneous
translation $(i,j)\mapsto(i+s,j+s)$. For $L,D$ this follows from
ratios of the $x_i$ and periodicity of the $\beta_i$.
For $U$ and $U^{-1}$, the factor $e^{2a\rho t h}$ acquired by a degree
$h$ symmetric polynomial cancels the factor $e^{-2a\rho^2s h}$ in
$v_{j+s}/v_{i+s}$, because $t=\rho s$.

For any such periodic, absolutely row-summable matrix $A$, define
\[
 (\Pi A)_{ij}=\sum_{n\in\Z} A_{i,j+ns}.
\]
Absolute convergence, decomposition of an intermediate index modulo
$s$, and simultaneous periodicity imply
$\Pi(AB)=(\Pi A)(\Pi B)$. Also
$\norm{\Pi A}\le\norm A$ and $\Pi I=I_s$.
Since $\Pi G=N$, periodizing the inverse proves the asserted bound for
$N^{-1}$.

Finally $c\ge\prod_{m\ge1}(1-e^{-2m})>5/6$.
At $a=1$, concavity gives
$1-e^{-\rho\theta}\ge1-e^{-\theta}\ge(1-e^{-1})\theta>3\theta/5$,
and $e^{1/2}<2$; substitution gives the stated constant $16$.
\end{proof}

\begin{corollary}[The width may be fixed]
For the resampling convention $a=\pi\alpha^2$, one may take the fixed
integer $\alpha=2$ in the inverse-conditioning and normalization
argument. If all $d$ gaps satisfy $\theta_i\ge d^{-C}$ for a fixed $C$,
the total logarithmic normalization loss is $O(d\log d)$.
\end{corollary}
\begin{proof}
The theorem bounds each inverse by a fixed constant times $d^{2C}$.
The separate diagonal has norm at most $e^{a/4}$, a fixed constant.
Taking powers of two large enough to normalize these factors costs
$O(\log d)$ bits per axis; logarithms add under tensor products.
\end{proof}

\paragraph{Scope and provenance.}
This strengthens the conditioning-only result in the earlier note
\emph{A Polynomial Condition Bound for Gaussian Resampling}. The
matrix is the one in Section~7 of the
\href{https://github.com/CrocSwap/integer-mult-bounds/tree/e6f9716621f4e01fc190617934cdfaf072bd976f}{inspected repository snapshot}.
Classical Newton interpolation and symmetric-polynomial identities are
used explicitly above; no generic-rank assumption is used. Two independent
agents checked the limiting and periodization arguments. Exact rational
controls cover negative and zero origins, finite $LDU$, inverse entries,
origin shifts and simultaneous periodicity.

This theorem does not supply the new numerical implementation at these
parameters. Stable factor application, in-place tensor composition,
coefficient generation and all tape movement must still be priced.
The repository's inspected conditional exponent remains
$\kappa=0.001672660197076918$; no larger final exponent is certified here.
\end{document}
